% Main figures: match filenames and labels to the manuscript.
% The source counts below match the supplied report. Verify against final export.
\newcommand{\SpectralFigureCaption}{%
\textbf{Width dependence of the posterior spectral statistic and domain coverage.}
Here $T=\|W_2\|_{\mathrm{op}}/\sqrt m$, the hidden-layer prior center is zero,
and $\sigma=1$, so the theorem permits any fixed cutoff
$a>a_*=r_0+2\sigma=2$. (a) Median, 95th, and 99th posterior percentiles of $T$
versus width. Points summarize three dataset/prior-center replicates by their
median, with whiskers showing their range. (b) Empirical coverage
$\widehat\pi(G_a)$ versus cutoff $a$, with one curve per width; curves and
shading show the pointwise median and range across replicates. Both panels
use unrestricted posterior samples and mark the threshold $a_*=2$.
For every fixed $a>2$, the theorem predicts coverage tending to one as width
increases; it makes no such assertion at $a=2$. Exact sample counts above 2
are reported in the accompanying table. Here $n=128$ and each target uses
four chains.}

\newcommand{\LossAcfFigureCaption}{%
\textbf{Width improves loss decorrelation, with diminishing returns in shallow networks.}
Autocorrelation of the summed training cross-entropy $V$ is shown over the
first 200 sampler lags for shallow ($L=2$) and deep ($L=3$) networks.
All runs use the same Metropolis-adjusted Gaussian-preserving Langevin
kernel with step size $h=0.01$ and $n=128$. Each of four chains contributes
its last 102,400 retained transitions, including repeats after rejection.
Curves are the pointwise median across three dataset/prior-center replicates
after averaging the four chain autocorrelations within each replicate;
shading shows the replicate range. The deep curves show a pronounced width
benefit, while the shallow improvement largely levels off after $m=256$.
The comparison concerns decorrelation per sampler iteration.}

% Supplementary templates: add actual findings after the analyses are complete.
% Change count/scope wording if the documented fallback is used.
\newcommand{\PredictionAcfFigureCaption}{%
\textbf{Prediction autocorrelation across widths.}
We compute autocorrelation separately for class-1 probabilities at eight
fixed held-out inputs, average over four chains and the eight inputs within
each replicate, and show the median and range across three
dataset/prior-center replicates. Each chain contributes the same final
102,400 retained transitions used in the loss analysis, with $h=0.01$.
Colors identify width. This tests whether the width pattern observed for
training loss also appears in predictive observables.}

\newcommand{\StepSizeFigureCaption}{%
\textbf{Step-size sensitivity of the deep width comparison.}
Training-loss autocorrelation is compared at widths 32 and 256 for two
step sizes of the same adjusted Langevin kernel. The horizontal axis is
the discrete algorithmic lag $kh$. Each chain contributes the final 512
units of retained algorithmic duration; curves show the median and range
of four-chain averages across three dataset/prior-center replicates.
The comparison tests whether the endpoint width ordering persists when
the step size is halved.}
